\chapter{Verwandte Arbeiten}


% Hilfsdefinition für linksbündige Spalten mit fester Breite und Umbruch
\newcolumntype{L}[1]{>{\raggedright\arraybackslash}p{#1}}

% Start der xltabular-Umgebung
% {\textwidth} = nutzt die volle Seitenbreite
% L{4.5cm} = erste Spalte fix 4.5cm, linksbündig
% X = mittlere Spalte nutzt den restlichen Platz flexibel
% L{3.5cm} = letzte Spalte fix 3.5cm, linksbündig
\begin{xltabular}{\textwidth}{L{4.5cm} X L{3.5cm}}
    
    \caption{Übersicht von verwandten Arbeiten, sortiert nach Erscheinungsjahr} \label{tab:literaturuebersicht} \\
    
    \toprule
    \textbf{Titel (Autoren)} & \textbf{Relevanz für die Arbeit} & \textbf{Geeignet für Kapitel} \\
    \midrule
    \endfirsthead % Kopfzeile nur auf der ersten Seite
    
    \multicolumn{3}{c}{\textit{Fortsetzung von der vorherigen Seite}} \\
    \toprule
    \textbf{Titel (Autoren)} & \textbf{Relevanz für die Arbeit} & \textbf{Geeignet für Kapitel} \\
    \midrule
    \endhead % Kopfzeile auf allen Folgeseiten
    
    \bottomrule
    \multicolumn{3}{r}{\textit{Weiter auf der nächsten Seite}} \\
    \endfoot % Fußzeile auf allen umgebrochenen Seiten
    
    \bottomrule
    \endlastfoot % Fußzeile nur am Ende der gesamten Tabelle

    % --- TABELLENINHALT START ---
    
    A Comprehensive Review on File Containers-Based Image and Video Forensics (Yang u.\,a. 2025)\cite{yang_comprehensive_2025} & 
    Aktueller Überblick über Container-Forensik & 
    Container-Forensik \\ \midrule
    
    Media source similarity hashing (MSSH): A practical method for large-scale media investigations (Klier und Baier 2025)\cite{klier_media_2025} & 
    Similarity Hashing für JPEG samt GitHub & 
    Similarity Hashing \\ \midrule
    
    Learning Hierarchical Fingerprints via Multi-Level Fusion for Video Integrity and Source Analysis (Y. Li u.\,a. 2024)\cite{li_learning_2024} & 
    Argumentiert, dass Container-Spuren durch Re-Muxing leicht manipulierbar sind & 
    Container-Forensik \\ \midrule
    
    Forensic Analysis of Binary Structures of Video Files (Hasan u.\,a. 2021)\cite{hasan_forensic_2021} & 
    Untersucht binäre Unterschiede in ftyp- und moov-Boxen zwischen Kameraherstellern & 
    Aufbau MP4, Container-Forensik \\ \midrule
    
    Forensic Analysis of Video Files Using Metadata (Xiang u.\,a. 2021)\cite{xiang_forensic_2021} & 
    Nutzt Vektorisierung von Metadaten-Bäumen für MP4 & 
    Aufbau MP4, Container-Forensik \\ \midrule
    
    Digital Video Source Identification Based on Container's Structure Analysis (Ramos López u.\,a. 2020)\cite{ramos_lopez_digital_2020} & 
    Nutzt hierarchische Container-Strukturen diverser Video-Dateien zur unüberwachten Quellenerkennung anhand der VISION- und ACID-Datensätze & 
    Aufbau MP4, Container-Forensik \\ \midrule
    
    Efficient Video Integrity Analysis Through Container Characterization (Yang u.\,a. 2020)\cite{yang_efficient_2020} & 
    Etabliert einen quantitativen Ansatz mittels Bag-of-Words-Modell und Entscheidungsbäumen zur Integritätsprüfung. Liefert zudem den EVA-7K-Datensatz (manipulierte und komprimierte Social-Media-Videos auf Basis „VISION“ von Shullani u.\,a.) & 
    Container-Forensik, Video-Datensätze \\ \midrule
    
    Information technology - Coding of audio-visual objects - Part 14: MP4 file format (ISO 2020)\cite{iso_14496-14_2020} & 
    ISO-Norm, definiert die formale Spezifikation und hierarchische Objektstruktur (Atome/Boxen) des MP4-Formats & 
    Aufbau MP4 \\ \midrule
    
    A Video Forensic Framework for the Unsupervised Analysis of MP4-Like File Container (Iuliani u.\,a. 2019)\cite{iuliani_video_2019} & 
    Formalisiert MP4 als Baumstruktur & 
    Aufbau MP4, Container-Forensik \\ \midrule
    
    VISION: a video and image dataset for source identification (Shullani u.\,a. 2017)\cite{shullani_vision_2017} & 
    Liefert Datensatz für native Videoquellen (Dataset „VISION“), Grundlage für Social-Media-Variante (Dataset „EVA-7K“ von Yang u.\,a.) & 
    Video-Datensätze \\ \midrule
    
    Forensic analysis of video file formats (Gloe u.\,a. 2014)\cite{gloe_forensic_2014} & 
    Zeigen auf, dass Hersteller und Software-Encoder spezifische Spuren in der Struktur von AVI- und MP4-Dateien hinterlassen, etwa durch die Anordnung und Benennung von Datenblöcken & 
    Aufbau MP4, Container-Forensik \\
    
    % --- TABELLENINHALT ENDE ---

\end{xltabular}



%%%%%

\section{Container-Forensik für Video- und Bilddaten}

\begin{table}[H]
    \centering
\caption{Verfahrensübersicht zur Container-Forensik}
\label{tab:containerforensik}
    \begin{tabular}{>{\raggedright\arraybackslash}p{0.17\linewidth}>{\raggedright\arraybackslash}p{0.1\linewidth}>{\raggedright\arraybackslash}p{0.5\linewidth}>{\raggedright\arraybackslash}p{0.1\linewidth}}\toprule
         \textbf{Verfahren}&  \textbf{Container}&  \textbf{Beschreibung}& \textbf{Quelle}\\\midrule
         n-Gramme&  jpeg&  Lineare Sequenzierung der Datei-Marker in Sets, flache Linearisierung& Klier und Baier (2025) \cite{klier_media_2025}\\\hline
 Bag-of-Words (BoW)& mp4& Zählt ausschließlich die absolute Häufigkeit, mit der ein bestimmtes Container-Symbol im MP4-Baum auftritt, und überführt diese Zählwerte in einen Feature-Vektor&Yang u. a. (2020)\\\hline
 Type-Aware (Datentyp-spezifisch)& mp4, mov& Unterscheidet Container-Symbole strikt nach Datentyp, um die Dimensionalität des Vektors zu reduzieren,  Vorhandensein einer bestimmten Box werden wie bei BoW als Häufigkeiten gezählt und kontinuierliche Felder (Bitrate, Dauer, Videoauflösung) werden als ihre tatsächlichen numerischen Werte werden direkt in den Feature-Vektor kopiert (Ist dies noch Container-Forensik?)&Xiang u.a. (2021)\\\hline
 Hierarchische Strukturanalyse& mp4, mov, 3gp& Extrahiert Container-Atome in vier starre Wörterbuch-Repräsentationen (PathField, PathFieldValue, PathOrderField, PathOrderFieldValue)&Ramos López u.a. (2020)\\ \hline
 Analyse binärer Strukturen& mp4, mov, avi& Manueller Hexadezimal-Abgleich von Box-Größen und Metadaten-Tags&Hasan u. a. (2021)\\ \bottomrule
    \end{tabular}
\end{table}

\subsection{State of the Art der Container-Forensik}

\begin{itemize}
    \item Containerforensik für Video/Bilder ist weiterhin noch relevant, wie ein systematischer Review-Artikel basierend auf 53 Publikationen von Yang u.a (2025) \cite{yang_comprehensive_2025} aufzeigt
    \begin{itemize}
        \item State-of-the-Art-Überblick zu Container-Forensik für Bilder und Videos
        \item Artikel erhebt keine eigenen Primärdaten, sondern aggregiert und taxiert bestehende Methoden
        \item Liefert Systematik und die Vergleichsmetriken
        \item 6-stufige Pipeline für die quantitative Analyse von Videodaten
        \item Feature-Repräsentation ausschließlich Bag-of-Words- und Type-Aware-Modelle für Video-Container  \cite[Kap. 4.3.4]{yang_comprehensive_2025}
        \item \textbf{Erkenntnisse}: Erhalt von Pfad und relativer Reihenfolge identischer Boxen als derzeit robustester Ansatz identifiziert
    \end{itemize}
\end{itemize}

\subsection{Containerstruktur und forensischer Nutzen}

Die Containerstruktur trägt forensisch nutzbare Informationen:

\begin{enumerate}
    \item Gloe u. a. (2014)\cite{gloe_forensic_2014}: Fundament der strukturellen und binären Analyse
    \item Iuliani u. a. (2019)\cite{iuliani_video_2019}: Basis der Hierarchischen Strukturanalyse\end{enumerate}

\subsection{Verfahren}

Bestehende Verfahren der Container-Forensik umgehen die rechenintensive Videokodierung, indem sie die hierarchische Struktur und Metadaten des MP4-Formats analysieren. Die hier analysierte Literatur weist hierbei vier methodische Richtungen auf, und zwar Analyse binärer Strukturen, Metadaten-Vektorisierung, hierarchische Strukturanalyse und quantifizierbare Container-Charakterisierung:

\begin{enumerate}
    \item Ramos López u. a. (2020)\cite{ramos_lopez_digital_2020}: Hierarchische Strukturanalyse (Pfad, Pfad + Wert, Pfad + Reihenfolge, Pfad + Reihenfolge + Wert)
    \item Hasan u. a. (2021)\cite{hasan_forensic_2021}: Analyse binärer Strukturen (Manueller und visueller Hex-Abgleich von Containerstrukturen)
    \begin{enumerate}
        \item Keine Erwähnung in Yangs Übersichtswerk (unzureichend belegter Ansatz?)
    \end{enumerate}
    \item Xiang u. a. (2021)\cite{xiang_forensic_2021}: Metadaten-Vektorisierung (Baumstruktur wird in ein Set aus Strings übersetzt und vektorisiert)
    \item Yang u. a. (2020)\cite{yang_efficient_2020}: Anwendung quantitativer Ansätze mittels Bag-of-Words-Modellen und Entscheidungsbäumen zur Integritätsprüfung
\end{enumerate}


\textbf{Kein Fokus} der Arbeit, aber als Reminder: Zuverlässigkeit von Containerstrukturen?

\begin{itemize}
    \item Y. Li u. a. (2024)\cite{li_learning_2024}: Methodologische Kritik (Re-Muxing-Vulnerabilität)
\end{itemize}



%%%%%

\section{Similarity Hashing in der digitalen Forensik}



%%%%%

\section{Video-Datensätze}

\begin{enumerate}
    \item VISION (Shullani u. a. 2017)\cite{shullani_vision_2017}: Videos Smartphone, Social Media (Facebook, YouTube, Whattsapp) \url{https://lesc.dinfo.unifi.it/VISION/}
    \item EVA-7K (Yang u. a. 2020)\cite{yang_efficient_2020}: Videos Social Media (Facebook, TikTok, Weibo, YouTube), Software-Tools (Avidemux, Exiftool, ffmpeg, Kdenlive, Adobe Premiere) \url{https://drive.google.com/drive/folders/1_J04iEbhZTxQgJC7cE3fSG8vkP1by8i1}
\end{enumerate}